\documentclass[12pt]{article}

% a template that a friend gave, it's worked well enough for me
% i have added some packages and stuff that have proved useful

\usepackage{fancyhdr}
\usepackage{tipa}
\usepackage{fontspec}
\usepackage{amsfonts}
\usepackage{enumitem}
\usepackage[margin=1in]{geometry}
\usepackage{graphicx}
\usepackage{float}
\usepackage{amsmath}
\usepackage{braket}
\usepackage{amssymb}
\usepackage{booktabs}
\usepackage{hyperref}
\usepackage{mathtools}
\usepackage{xcolor}
\usepackage{float}
\usepackage{algpseudocodex}
\usepackage{titlesec}
\usepackage{bbm}

\pagestyle{fancy}
\fancyhf{} % sets both header and footer to nothing
\lhead{Kevin Sheng}
\setmainfont{Comic Neue}
\renewcommand{\headrulewidth}{1pt}
\setlength{\headheight}{0.75in}
\setlength{\oddsidemargin}{0in}
\setlength{\evensidemargin}{0in}
\setlength{\voffset}{-.5in}
\setlength{\headsep}{10pt}
\setlength{\textwidth}{6.5in}
\setlength{\headwidth}{6.5in}
\setlength{\textheight}{8in}
\renewcommand{\headrulewidth}{0.5pt}
\renewcommand{\footrulewidth}{0.3pt}
\setlength{\textwidth}{6.5in}
\usepackage{setspace}
\usepackage{multicol}
\usepackage{float}
\setlength{\columnsep}{1cm}
\setlength\parindent{24pt}
\usepackage [english]{babel}
\usepackage [autostyle, english = american]{csquotes}
\MakeOuterQuote{"}

\setlength{\parskip}{6pt}
\setlength{\parindent}{0pt}

\titlespacing\section{0pt}{12pt plus 4pt minus 2pt}{0pt plus 2pt minus 2pt}
\titlespacing\subsection{0pt}{12pt plus 4pt minus 2pt}{0pt plus 2pt minus 2pt}
\titlespacing\subsubsection{0pt}{12pt plus 4pt minus 2pt}{0pt plus 2pt minus 2pt}

\hypersetup{colorlinks=true, urlcolor=blue}

\newcommand{\correction}[1]{\textcolor{red}{#1}}


\rhead{Math 110A}

\makeatletter
\def\@seccntformat#1{%
  \expandafter\ifx\csname c@#1\endcsname\c@section\else
  \csname the#1\endcsname\quad
  \fi}
\makeatother

\newcommand{\lra}{\xLeftrightarrow}
\newcommand{\ra}{\xRightarrow}

\begin{document}

\section{Problem 1}

\subsection{Zero Commutes}

For $r \cdot 0$, the proof proceeds like so:
\begin{align*}
  0 + 0 = 0
   & \implies r \cdot (0 + 0) = r \cdot 0       \\
   & \implies r \cdot 0 + r \cdot 0 = r \cdot 0
\end{align*}
Since $f(x)=x+r \cdot 0$ is a bijection from $R$ to $R$,
$0$ can be the only element that leaves $r \cdot 0$ unchanged when
added to it.
Thus, $r \cdot 0 = 0$.

For the other multiplication, the proof proceeds the same except for us
having to place the $r$ to the right of $0+0$ instead of the left. $\square$

\subsection{Negation Too}

For $r \cdot (-1) = -r$, notice that
\begin{align*}
  1+(-1)=0
   & \implies r \cdot (1+(-1))=r \cdot 0=0 \\
   & \implies r \cdot 1 + r \cdot (-1) = 0 \\
   & \implies r + r \cdot (-1) = 0
\end{align*}
Since $-r$ is defined as the unique element that makes $0$ when added to $r$, $r \cdot (-1) = -r$.


If we wanted to do this for $(-1) \cdot r$, it's the same except we have to
multiply by putting $r$ to the right of $1+(-1)$. $\square$

\section{Problem 2}

$\Z^{\ge 0}$ does not form a ring.

Consider the function $f(x)=x+1$.
This does not make a bijection from $\Z^{\ge 0}$ onto itself, because
no number gives $0$ when plugged into it.

\pagebreak

\section{Problem 3 and 4}

I'll show a more general version, which is that $R\braket{x_i, i \in I}$ is a ring
and that all of its units have to be constants.
$R[x]$ and $R\braket{x, y}$ are just the above set with $I=\{x\}$ and $I=\{x, y\}$
respectively, so this solves both problems at once.

I'm not sure if this was defined accurately in class,
so I'll express an element $r$ of this ring in the form $\sum_{\phi} r_\phi x^\phi$.
The idea is that $r$ is a mapping from $\hom_{set}([n], I)$ to $R$
that only gives nonzero coefficients on finitely many inputs.
In other words, I'm both using $r$ as notation for the polynomial function itself
as well as identification for its coefficients.
Also, here the summation is over all $\phi \in \hom_{\text{set}}([n], I)$.

\subsection{Addition Commutes}

I'll be using $RI$ as an abbreviation for $R\braket{x_i, i \in I}$.

Consider any two $r, s \in RI$:
\begin{align*}
  r+s
  &= \sum_{\phi} \left(r_\phi+s_\phi\right)x^\phi \\
  &= \sum_{\phi} \left(s_\phi+r_\phi\right)x^\phi \\
  &= s+r
\end{align*}

\subsection{Addition Associates}

oh my god bruh
\begin{align*}
  r+(s+t)
  &= \sum_{\phi} r_\phi x^\phi+\sum_{\phi} \left(s_\phi+t_\phi\right)x^\phi \\
  &= \sum_{\phi} \left(r_\phi+(s_\phi+t_\phi)\right)x^\phi \\
  &= \sum_{\phi} \left((r_\phi+s_\phi)+t_\phi\right)x^\phi \\
  &= \sum_{\phi} (r_\phi+s_\phi) x^\phi+\sum_{\phi} t_\phi x^\phi \\
  &=(r+s)+t
\end{align*}

\subsection{Addition is a Bijection}

Fix an $r$ and consider $a_r(s): r+s$.

For \textbf{injectivity}, consider $s, t \in RI$ with $s \ne t$.
This means $\exists \psi \in \hom_{\text{set}}([n], I): s_\psi \ne t_\psi$.

We then consider $a_r(s)$ and $a_r(t)$:
\begin{align*}
  a_r(s)=\sum_{\phi} (r_\phi+s_\phi)x^\phi &&
  a_r(t)=\sum_{\phi} (r_\phi+t_\phi)x^\phi
\end{align*}
Since the summation is only one level deep, $a_r(s)_\psi=r_\psi+s_\psi$ and $a_r(t)_\psi=r_\psi+t_\psi$.
By the inequality established earlier, this implies that $a_r(s) \ne a_r(t)$.

For \textbf{surjectivity}, fix any $t \in RI$.
It remains to find an $s \in RI$ s.t. $a_r(s)=t$.

I propose that $s=\sum_{\phi} (t_\phi-r_\phi)x^\phi$ is our answer.
$r_\phi \in R$, so we can negate elements without any issues:
\begin{align*}
  a_r(s)
  &= \sum_{\phi} r_\phi x^\phi + \sum_{\phi} (t_\phi-r_\phi)x^\phi \\
  &= \sum_{\phi} (r_\phi+t_\phi-r_\phi)x^\phi \\
  &= \sum_{\phi} t_\phi x^\phi \\
  &= t
\end{align*}

\subsection{Multiplication Associates}

I have no idea what the third letter associated with $\phi$ and $\psi$ is.
I'm just gonna use $\rho$, please don't kill me.
\begin{align*}
  r \cdot (s \cdot t)
  &= \sum_{\rho} r_\rho x^\rho \cdot \sum_{\phi} \sum_{\psi} s_\phi t_\psi x^{\phi \sqcup \psi} \\
  &= \sum_{\rho} \sum_{\phi} \sum_{\psi} r_\rho (s_\phi t_\psi) x^{\rho \sqcup \phi \sqcup \psi} \\
  &= \sum_{\rho} \sum_{\phi} \sum_{\psi} (r_\rho s_\phi) t_\psi x^{\rho \sqcup \phi \sqcup \psi} \\
  &= \sum_{\rho} \sum_{\phi} r_\rho s_\phi x^{\rho \sqcup \phi} \cdot \sum_{\psi} t_\psi x^\psi \\
  &= (r \cdot s) \cdot t
\end{align*}

\subsection{Distributivity}

For left distributivity, we have
\begin{align*}
  r \cdot (s + t)
  &= \sum_{\phi} r_\phi x^\phi \cdot \sum_{\phi} (s_\phi+t_\phi)x^\phi \\
  &= \sum_{\phi} \sum_{\psi} r_\phi(s_\psi+t_\psi) x^{\phi \sqcup \psi} \\
  &= \sum_{\phi} \sum_{\psi} (r_\phi s_\psi + r_\phi t_\psi) x^{\phi \sqcup \psi} \\
  &= \sum_{\phi} \sum_{\psi} r_\phi s_\psi x^{\phi \sqcup \psi} + \sum_{\phi} \sum_{\psi} r_\phi t_\psi x^{\phi \sqcup \psi} \\
  &= r \cdot s + r \cdot t
\end{align*}
since the coefficients are in $R$ and we can apply distributivity to them.

The proof for right distributivity is basically the same:
\begin{align*}
  (s + t) \cdot r
  &= \sum_{\phi} (s_\phi+t_\phi)x^\phi \cdot \sum_{\phi} r_\phi x^\phi \\
  &= \sum_{\phi} \sum_{\psi} (s_\psi+t_\psi)r_\phi x^{\phi \sqcup \psi} \\
  &= \sum_{\phi} \sum_{\psi} (s_\psi r_\phi + t_\psi r_\phi) x^{\phi \sqcup \psi} \\
  &= \sum_{\phi} \sum_{\psi} s_\phi r_\psi x^{\phi \sqcup \psi} + \sum_{\phi} \sum_{\psi} t_\phi r_\psi x^{\phi \sqcup \psi} \\
  &= s \cdot r + t \cdot r
\end{align*}

\subsection{Zero}

I propose that $0$ in the ring, or an $r_\phi$ that gives $0$ on every $\phi$, is our additive identity:
\begin{align*}
  r+0
  &= \sum_{\phi} r_\phi x^\phi + \sum_{\phi} 0 \cdot x^\phi \\
  &= \sum_{\phi} (r_\phi+0) x^\phi \\
  &= r
\end{align*}

\subsection{One}

Our multiplicative identity should be an $r_\phi$ that gives $0$ on every $\phi$
except on the single $\psi: [0] \to I$, for which it gives $1$.
Since $\phi \sqcup \psi = \phi$ for any $\phi \in \hom_{\text{set}}([n], I)$:
\begin{align*}
  r \cdot 1
  &= \sum_{\phi} r_\phi x^\phi \cdot (1 \cdot x^\psi) \\
  &= \sum_{\phi} (r_\phi \cdot 1)x^{\phi \sqcup \psi} \\
  &= \sum_{\phi} r_\phi x^\phi \\
  &= r
\end{align*}

Because I'm paranoid, I'll also prove the same thing for
$1 \cdot r$ even though it's basically the same thing:
\begin{align*}
  1 \cdot r
  &= (1 \cdot x^\psi) \cdot \sum_{\phi} r_\phi x^\phi \\
  &= \sum_{\phi} (1 \cdot r_\phi)x^{\psi \sqcup \phi} \\
  &= \sum_{\phi} r_\phi x^\phi \\
  &= r
\end{align*}

\pagebreak

\section{Problem 5}

\subsection{\texorpdfstring{$\mathrm{End}(V)$}{End(V)} is a Ring}

It was proven in 115A that addition and composition of linear transformations
still results in another linear transformation, so I will not be proving
that these operations are well-defined.

\subsubsection{Addition Commutes}

Since we know vector space elements commute,
\begin{align*}
  (\phi+\psi)(v)
   & = \phi(v)+\psi(v) \\
   & = \psi(v)+\phi(v) \\
   & = (\psi+\phi)(v)
\end{align*}

\subsubsection{Addition Associates}

I'll just use the alphabet here.
Using $\rho$ as a counterpart for $\phi$ and $\psi$ felt a little weird:
\begin{align*}
  (a+(b+c))(v)
   & = a(v)+(b+c)(v)    \\
   & = a(v)+(b(v)+c(v)) \\
   & = (a(v)+b(v))+c(v) \\
   & = (a+b)(v)+c(v)    \\
   & = ((a+b)+c)(v)
\end{align*}

\subsubsection{Addition is a Bijection}

Fix an $a \in \mathrm{End}(v)$ and consider $f_a(\phi)=(a+\phi)$.

To show \textbf{injectivity}, consider $b, c \in \mathrm{End}(V)$.
If $b \ne c$, $\exists v \in V: b(v) \ne c(v)$.

We also have $(f_a(b))(v)=a(v)+b(v)$ and $(f_a(c))(v)=a(v)+c(v)$.
Since $b(v) \ne c(v)$, then $a(v)+b(v) \ne a(v)+c(v)$
and therefore $f_a(b) \ne f_a(c)$.

For \textbf{surjectivity}, take any $b \in \mathrm{End}(v)$.

I propose that when we add $a$ to $c(v)=b(v)-a(v)$, where subtraction
means the addition of the vector space's additive inverse, we will obtain $b$.

The proof of this is straightforward:
\begin{align*}
  (c+a)(v)
   & = c(v)+a(v)      \\
   & = b(v)-a(v)+a(v) \\
   & = b(v)
\end{align*}

\subsubsection{Multiplication Associates}

In general, function composition is associative:
\begin{align*}
  (f \circ (g \circ h))(v)
   & = f((g \circ h)(v))        \\
   & = f(g(h(v)))               \\
   & = (f \circ g)(h(v))        \\
   & = ((f \circ g) \circ h)(v)
\end{align*}
so "multiplication" in this ring associates.

\subsubsection{Distributivity}

Using the properties of linear transformations, we have for left distributivity:
\begin{align*}
  (a \cdot (b+c))(v)
   & = a((b+c)(v))                     \\
   & = a(b(v)+c(v))                    \\
   & = a(b(v))+a(c(v))                 \\
   & = (a \cdot b)(v) + (a \cdot c)(v) \\
   & = ((a \cdot b) + (a \cdot c))(v)
\end{align*}
and for right distributivity:
\begin{align*}
  ((a+b) \cdot c)(v)
   & = (a+b)(c(v))                     \\
   & = a(c(v))+b(c(v))                 \\
   & = (a \cdot c)(v) + (b \cdot c)(v) \\
   & = ((a \cdot c) + (b \cdot c))(v)
\end{align*}

\subsubsection{Zero}

Our zero element should be $\phi(v)=0$, where $0$ is the zero element of the vector space.

This gives us
\begin{align*}
  (\phi + \psi)(v) & = \phi(v)+\psi(v) & (\psi + \phi)(v) & = \psi(v)+\phi(v) \\
                   & = 0+\psi(v)       &                  & = \psi(v)+0       \\
                   & = \psi(v)         &                  & = \psi(v)
\end{align*}

\subsubsection{One}

Our identity element should be the identity transformation $\phi(v)=v$.

This way, we have both
\begin{align*}
  (\phi \cdot \psi)(v) & = \phi(\psi(v)) & (\psi \cdot \phi)(v) & = \psi(\phi(v)) \\
                       & = \psi(v)       &                      & = \psi(v)
\end{align*}

\subsection{A Vector Space Too}

Since addition forms a bijection, this implies the existence of an
additive inverse for $\mathrm{End}(V)$.
There's just some more axioms.

\subsubsection{Field Multiplication Commutes}

We have
\begin{align*}
  ((\lambda_1 \lambda_2) \cdot \phi)(v)
   & = \phi(\lambda_1\lambda_2 \cdot v)            \\
   & = \phi(\lambda_1(\lambda_2 \cdot v))          \\
   & = (\lambda_1 \cdot (\lambda_2 \cdot \phi))(v)
\end{align*}

\subsubsection{Distributivity}

We have
\begin{align*}
  (\lambda \cdot (\phi + \psi))(v)
   & = (\phi+\psi)(\lambda \cdot v)                      \\
   & = \phi(\lambda \cdot v) + \psi(\lambda \cdot v)     \\
   & = (\lambda \cdot \phi)(v) + (\lambda \cdot \psi)(v)
\end{align*}
and
\begin{align*}
  ((\lambda_1+\lambda_2) \cdot \phi)(v)
   & = \phi((\lambda_1+\lambda_2)v)                          \\
   & = \phi(\lambda_1 \cdot v + \lambda_2 \cdot v)           \\
   & = \phi(\lambda_1 \cdot v) + \phi(\lambda_2 \cdot v)     \\
   & = (\lambda_1 \cdot \phi)(v) + (\lambda_2 \cdot \phi)(v)
\end{align*}

\pagebreak

\section{Problem 6}

\subsection{Part A}

For $x$, we have
\begin{align*}
  x(\lambda f + g)
   & = x \cdot (\lambda f + g)        \\
   & = x \cdot \lambda f + x \cdot g        \\
   & = \lambda(x \cdot f) + x \cdot g \\
   & = \lambda \cdot x(f) + x(g)
\end{align*}
and for $\partial_x$, we have
\begin{align*}
  \partial_x(\lambda f + g)
   & = \frac{d}{dx} (\lambda f + g)                                          \\
   & = \frac{d}{dx} \lambda f + \frac{d}{dx} g     & \text{By Addition Rule} \\
   & = \lambda \frac{df}{dx} + \frac{dg}{dx}       & \text{By Product Rule}  \\
   & = \lambda \cdot \partial_x(f) + \partial_x(g)
\end{align*}

\subsection{Part B}

When we multiply then differentiate, we get
\begin{align*}
  \frac{d}{dx} x\sum_{i=0}^{n} r_i x^i
  &= \frac{d}{dx} \sum_{i=0}^{n} r_i x^{i+1} \\
  &= \sum_{i=0}^{n} r_i(i+1) x^i
\end{align*}
OTOH, when we switch the order, we get
\begin{align*}
  x\frac{d}{dx} \sum_{i=0}^{n} r_i x^i
  &= x\sum_{i=0}^{n} r_i \cdot i \cdot x^{i-1} \\
  &= \sum_{i=0}^{n} r_i \cdot i \cdot x^i
\end{align*}

Subtracting the second from the first gives
\[\sum_{i=0}^{n} r_i(i+1)x^i - \sum_{i=0}^{n} r_i \cdot i \cdot x^i
=\boxed{\sum_{i=0}^{n} r_i x^i}\]

\end{document}
